\section{Related Work}
\label{sec:related}

\textbf{FANET routing.} Bekmezci~\emph{et~al.}~\cite{bekmezci2013fanet} define FANETs and point out what
separates them from mobile and vehicular ad hoc networks: higher speeds, 3-D motion and frequent topology
changes. Khan~\emph{et~al.}~\cite{khan2020survey} survey FANET routing schemes by family (topology-based,
position-based, hierarchical, learning-based) and identify link instability as a central open problem.
\Cref{tab:related} positions the protocols studied here.

\textbf{Topology-based routing.} DSDV~\cite{perkins1994dsdv} is proactive: every node keeps a route to every
destination, refreshed by periodic and triggered table updates with sequence numbers. Its control traffic
grows with the network, and its routes lag behind the topology between updates. GRAd~\cite{poor2001grad}
builds a cost-to-destination gradient on demand and forwards by broadcast: any neighbour that hears a packet
and is closer to the destination relays it. On UavNetSim-v1 GRAd disables MAC acknowledgements, so it has no
retransmissions and no single next hop.

\textbf{Position-based routing.} GPSR~\cite{karp2000gpsr} forwards greedily to the neighbour closest to the
destination and falls back on perimeter routing around voids. The platform's Greedy forwarding implements
the greedy mode only: when no neighbour makes progress, the packet waits in a waiting list. Greedy needs no
route state and adapts per packet, but it selects by distance alone and ignores how long the chosen link
will last.

\textbf{Prediction-based routing.} Su~\emph{et~al.}~\cite{su2001mobility} introduced the link expiration
time: given the positions and velocities of two nodes and their range, the time until they separate has a
closed form. OPAR~\cite{gharib2021opar} computes shortest paths constrained by such link lifetimes. On the
platform, OPAR is centralised: the source reads the true positions and velocities of all UAVs. It is
therefore an upper reference with global knowledge rather than a distributed competitor. TARRAQ~\cite{cui2022tarraq}
predicts link duration with a Kalman filter and feeds it into adaptive Q-learning.

\textbf{Learning-based routing.} Q-Routing~\cite{boyan1993qrouting} learns, for each neighbour, an estimate
of the delivery time through it. QGeo~\cite{jung2017qgeo} applies Q-learning to geographic forwarding in
unmanned robotic networks, and Q-FANET~\cite{dacosta2021qfanet} refines it for FANETs. Learning-based
protocols react to observed outcomes. They need feedback before they adapt, and they do not explain why a
packet was lost.

\textbf{Gap.} Judging prediction fairly requires a measurement of how many losses are link breaks, separated
from collisions and queueing, and a comparison that holds everything else fixed, including the MAC and the
reaction to failures. This paper provides both, with a distributed predictor that uses only information
carried in hello messages.

\begin{table*}[t]
\centering
\caption{Protocols compared in this work, as implemented on UavNetSim-v1~v1.0.1 (``retry same hop'': the MAC resends to the same neighbour, up to five attempts).}
\label{tab:related}
\footnotesize
\setlength{\tabcolsep}{3pt}
\begin{tabular}{@{}lllll@{}}
\toprule
Protocol & Family & Decision & Mobility prediction & After an ACK timeout \\
\midrule
DSDV~\cite{perkins1994dsdv} & topology, proactive & distributed & no & retry same hop \\
GRAd~\cite{poor2001grad} & topology, on demand & distributed & no & n/a (no ACK) \\
Greedy~\cite{karp2000gpsr} & position & distributed & no & retry same hop \\
OPAR~\cite{gharib2021opar} & prediction & centralised & LET, oracle & retry same hop \\
Q-Routing~\cite{boyan1993qrouting} & learning & distributed & no & retry same hop \\
QGeo~\cite{jung2017qgeo} & learning, position & distributed & no & penalise $Q$, retry \\
Q-FANET~\cite{dacosta2021qfanet} & learning, position & distributed & no & penalise $Q$, retry \\
\textbf{LEAP-GR} (ours) & position, prediction & distributed & LET, from hellos & see \cref{sec:leap} \\
\bottomrule
\end{tabular}
\end{table*}
